\documentclass[12pt]{article}
% Adaptado de C:/GitHub/Mestrado/Template/Template.tex.
\usepackage[T1]{fontenc}
\usepackage[utf8]{inputenc}
\usepackage[brazilian]{babel}
\usepackage{graphicx,booktabs,array,amsmath,amssymb,microtype}
\usepackage{geometry}
\usepackage{pdflscape,caption}
\geometry{a4paper,margin=2.5cm}
\usepackage{parskip,csquotes}
\interfootnotelinepenalty=10000
\usepackage[backend=biber,style=abnt]{biblatex}
\addbibresource{references.bib}
\usepackage{hyperref}
\hypersetup{colorlinks=true,urlcolor=blue,linkcolor=blue,citecolor=blue}
\newcommand{\instituicao}{UNIVERSIDADE FEDERAL DE SANTA CATARINA}
\newcommand{\departamento}{Programa de Pós-Graduação em Economia}
\newcommand{\curso}{Mestrado}
\newcommand{\disciplina}{Desigualdade, Diversidade e Redes}
\newcommand{\titulo}{Da produção à demanda}
\newcommand{\subtitulo}{Estrutura da rede intersetorial de atribuição de CO\textsubscript{2} no Brasil e implicações para a coordenação da mitigação}
\newcommand{\autor}{Gabriel Pincelli Cintra}
\newcommand{\local}{Florianópolis}
\newcommand{\dataentrega}{18 de outubro de 2026}
\newcommand{\imprimirtitulodois}{
\begin{center}
  {\Large\bfseries\MakeUppercase{\instituicao}\par}
  \vspace{0.2em}{\large \departamento\par}
  \vspace{0.2em}{\normalsize\textbf{\disciplina}\ (\curso)\par}
  \vspace{1.4em}{\LARGE\bfseries \titulo\par}
  \vspace{0.25em}{\large\itshape \subtitulo\par}
  \vspace{1.5em}
\end{center}
\noindent\begin{minipage}{0.48\textwidth}
  \raggedright\normalsize \autor
\end{minipage}\hfill
\begin{minipage}{0.48\textwidth}
  \raggedleft\normalsize \local, \dataentrega
\end{minipage}
\vspace{1.0em}\hrule\vspace{1.0em}}
\begin{document}
\imprimirtitulodois
\begin{abstract}
Este ensaio examina as relações intersetoriais que compõem as diferenças entre as emissões atribuídas pela produção e pela demanda no Brasil. Combina a Matriz de Insumo-Produto de 2015, com 67 atividades, e intensidades de CO\textsubscript{2} em uma rede dirigida e ponderada. As forças e o número efetivo de contrapartes caracterizam volume e concentração; o filtro de disparidade seleciona relações conforme sua importância local. A rede quase completa apresenta perfis distintos de distribuição dos pesos. Energia e construção ilustram os papéis de origem emissora e destino da atribuição. O backbone retém 13,40\% das ligações e 70,69\% do peso intersetorial, com coberturas locais desiguais. A análise informa relações a investigar para a coordenação da mitigação, mas as limitações dos coeficientes restringem as conclusões ao cenário exploratório.
\end{abstract}
\noindent\textbf{Palavras-chave:} atribuição de emissões; insumo-produto; redes dirigidas e ponderadas; mitigação.

\input{introducao.tex}
\widowpenalty=10000
\clubpenalty=10000
\displaywidowpenalty=10000
% Sequência inspirada em Huang et al. (2019).
\input{dados.tex}
\section{Metodologia}
\input{modelo_atribuicao.tex}
\subsection{Construção da rede dirigida e ponderada}
\label{sec:construcao-rede}

A rede tem 67 nós, um por atividade da MIP, identificados pelos códigos do IBGE. Cada nó representa uma origem emissora nas ligações de saída e o setor de um produto final nas ligações de entrada, conforme a representação de \textcite{huang2019}. As categorias de uso final não constituem nós separados.

Para estudar relações entre atividades distintas, define-se
\begin{equation}
W_{ij}=\begin{cases}
F_{ij}, & i\ne j,\\
0, & i=j.
\end{cases}
\label{eq:rede-sem-diagonal}
\end{equation}
Cada $W_{ij}>0$ define uma aresta $i\rightarrow j$, cujo peso, em Gg de CO\textsubscript{2}, indica a emissão em $i$ atribuída ao produto final de $j$. As duas direções de um par são distintas. Como os pesos incorporam requisitos diretos e indiretos, concatenar arestas não representa uma sequência de etapas produtivas.

A diagonal $F_{ii}$ reúne emissões de $i$ atribuídas à demanda final de seus próprios produtos, incluindo requisitos indiretos. Ela permanece em $F$ e nas contas $c$ e $u$, mas é retirada de $W$, cujo peso total corresponde apenas às atribuições entre atividades distintas.

A construção preserva todos os nós e as arestas positivas, sem corte por magnitude. São conferidos o alinhamento dos códigos e a existência de pesos finitos e não negativos. A seleção de arestas é tratada na seção~\ref{sec:backbone}.

\subsection{Volume e concentração das relações}
\label{sec:volume-concentracao}

O volume e a concentração são calculados em $W$, antes da filtragem, por meio das forças e da disparidade dos pesos \parencite{huang2019,serrano2009}.

\paragraph{Volume.}
As forças de saída e entrada são as somas dos pesos em cada direção:
\begin{equation}
\begin{aligned}
s_i^{\mathrm{out}}&=\sum_{j\ne i}W_{ij}=c_i-F_{ii},\\
s_i^{\mathrm{in}}&=\sum_{j\ne i}W_{ji}=u_i-F_{ii}.
\end{aligned}
\label{eq:forcas-externas}
\end{equation}
A força de saída mede as emissões próprias de $i$ atribuídas aos produtos finais de outras atividades; a de entrada mede as emissões de outras atividades atribuídas aos produtos finais de $i$. A diferença entre força de entrada e força de saída é $s_i^{\mathrm{in}}-s_i^{\mathrm{out}}=u_i-c_i$. A leitura separada dos volumes revela atribuições que se compensam no saldo; as contas completas e a diagonal situam sua participação no total.

\paragraph{Concentração.}
Para perfis com força positiva, define-se a participação de cada contraparte:
\begin{equation}
p_{ij}^{\mathrm{out}}=\frac{W_{ij}}{s_i^{\mathrm{out}}},\qquad
p_{ji}^{\mathrm{in}}=\frac{W_{ji}}{s_i^{\mathrm{in}}},\quad j\ne i.
\label{eq:participacoes-externas}
\end{equation}
O perfil de saída distribui o volume intersetorial da origem entre destinos; o de entrada descreve as origens que compõem a atribuição intersetorial ao destino. Cada perfil soma um. A disparidade e seu inverso são
\begin{equation}
\begin{aligned}
Y_i^{\mathrm{out}}&=\sum_{j\ne i}(p_{ij}^{\mathrm{out}})^2,&
Y_i^{\mathrm{in}}&=\sum_{j\ne i}(p_{ji}^{\mathrm{in}})^2,\\
D_{2,i}^{\mathrm{out}}&=1/Y_i^{\mathrm{out}},&
D_{2,i}^{\mathrm{in}}&=1/Y_i^{\mathrm{in}}.
\end{aligned}
\label{eq:disparidade}
\end{equation}
$Y$ aumenta com a concentração; $D_2$ expressa a mesma informação como número efetivo de contrapartes. Para $k$ pesos iguais, $D_2=k$; quando uma relação domina, aproxima-se de um. As medidas são adimensionais e ficam indefinidas quando a força é zero, sem excluir o nó.

Os pesos e as participações identificam as contrapartes, cuja identidade não é informada por $D_2$. O grau conta as arestas positivas de entrada ou saída; a densidade é o número de arestas dividido por $67\times66$. Esses indicadores descrevem a presença de relações, enquanto a disparidade considera a distribuição dos pesos.

\subsection{Extração e avaliação do backbone}
\label{sec:backbone}

O filtro de disparidade de \textcite{serrano2009}, aplicado por \textcite{huang2019}, seleciona ligações destacadas na composição local dos pesos, considerando atividades de diferentes escalas. O backbone resultante não altera as contas ambientais completas.

\paragraph{Critério de seleção.}
O modelo nulo distribui a força entre $k$ ligações dividindo o intervalo unitário por $k-1$ pontos uniformemente escolhidos. Para $k>1$, a probabilidade de uma parcela atingir ou superar $p$ é $(1-p)^{k-1}$. Para a aresta $i\rightarrow j$, as caudas nas perspectivas de saída e entrada são
\begin{equation}
\begin{aligned}
a_{ij}^{\mathrm{out}}&=\left(1-\frac{W_{ij}}{s_i^{\mathrm{out}}}\right)^{k_i^{\mathrm{out}}-1},\\
a_{ij}^{\mathrm{in}}&=\left(1-\frac{W_{ij}}{s_j^{\mathrm{in}}}\right)^{k_j^{\mathrm{in}}-1}.
\end{aligned}
\label{eq:caudas-filtro}
\end{equation}
Os graus $k$ e as forças são os de $W$. A aresta é selecionada se $a_{ij}^{\mathrm{out}}<\alpha$ \textbf{ou} $a_{ij}^{\mathrm{in}}<\alpha$, registrando-se a extremidade responsável pela seleção. A saída destaca um destino nas atribuições da origem; a entrada destaca uma origem na composição do destino.\footnote{As caudas são calculadas apenas para graus pertinentes maiores que um. Se uma extremidade tem grau um, vale o teste da outra; se ambas têm grau um, a aresta é preservada por convenção de manutenção da relação exclusiva, inclusive em díades isoladas. Essa extensão dos casos degenerados do suplemento de Serrano et al. é registrada separadamente da seleção estatística.}

Adota-se $\alpha=0{,}05$, com $0{,}01$ e $0{,}10$ para sensibilidade. O filtro é aplicado uma vez a $W$ para cada valor, sem atualização dos perfis. A matriz $B_\alpha$ preserva os pesos originais nas arestas selecionadas e zero nas demais, sem renormalização ou corte absoluto adicional.

\paragraph{Avaliação da seleção.}
Comparam-se a fração de arestas retidas e a cobertura dos pesos, globalmente e por atividade:
\begin{equation}
\begin{aligned}
r_\alpha&=\frac{\sum_{ij}(B_\alpha)_{ij}}{\sum_{ij}W_{ij}},\\
r_{\alpha,i}^{\mathrm{out}}&=\frac{\sum_j(B_\alpha)_{ij}}{s_i^{\mathrm{out}}},\qquad
r_{\alpha,i}^{\mathrm{in}}=\frac{\sum_j(B_\alpha)_{ji}}{s_i^{\mathrm{in}}}.
\end{aligned}
\label{eq:cobertura-backbone}
\end{equation}
As razões são definidas apenas para denominadores positivos. A cobertura refere-se à parcela intersetorial de $W$ e permite distinguir retenção agregada de perdas por atividade e direção. A conectividade é descrita pelos graus, isolados e componentes fracamente conectados, obtidos ao desconsiderar a direção das arestas.

A sensibilidade é examinada pelas mudanças nas arestas, coberturas e posições setoriais entre os três valores de $\alpha$. O aninhamento dos conjuntos decorre da regra e não constitui evidência independente de robustez. As caudas do modelo nulo não medem incerteza ambiental, e a seleção não estabelece prioridade de mitigação.

\section{Resultados}
\label{sec:resultados}

Os resultados correspondem ao cenário ambiental exploratório descrito na seção~\ref{sec:dados}.

\subsection{Contas ambientais, volume e concentração}

As duas contas totalizam 683,70 mil Gg de CO\textsubscript{2}, dos quais 392,92 mil Gg estão na diagonal de $F$ e 290,78 mil Gg nas relações entre atividades distintas em $W$. A igualdade entre os totais pela produção e pela demanda decorre da identidade contábil.

A rede completa tem 67 nós e 4\,290 arestas, com densidade de 97,01\%. Serviços domésticos (9700) é o único isolado: sua conta permanece na diagonal de $F$. As outras 66 atividades estão ligadas entre si nas duas direções, com graus iguais a 65. O número de ligações pouco distingue os perfis setoriais.

Os pesos, entretanto, distinguem volume e concentração (Figura~\ref{fig:indicadores-resultados}). Energia (3500) e água, esgoto e resíduos (3680) distribuem volumes externos próximos, de 31,32 e 31,23 mil Gg, mas apresentam 20,79 e 9,88 destinos efetivos. Assim, o volume de energia é mais distribuído entre contrapartes, enquanto o de água e resíduos é mais concentrado.

Energia e construção ocupam posições distintas nas contas completas (Tabela~\ref{tab:contas-energia-construcao}). Em energia, a produção supera a atribuição pela demanda; em construção ocorre o inverso. A pequena saída intersetorial de construção não indica baixa emissão própria: a maior parte de sua conta pela produção é atribuída à demanda final da própria atividade e permanece fora de $W$.

\clearpage
\begin{table}[!ht]
\centering
\caption{Contas e volumes intersetoriais de energia e construção.}
\label{tab:contas-energia-construcao}
\small
\begin{tabular}{@{}lrr@{}}
\toprule
Medida & Energia (3500) & Construção (4180) \\
\midrule
Produção, $c$ & 56,62 & 38,84 \\
Demanda, $u$ & 28,23 & 59,40 \\
Atribuição à própria atividade, $F_{ii}$ & 25,30 & 35,91 \\
Saída externa, $s^{\mathrm{out}}$ & 31,32 & 2,93 \\
Entrada externa, $s^{\mathrm{in}}$ & 2,92 & 23,48 \\
Saldo, $u-c$ & $-28,39$ & 20,56 \\
\bottomrule
\end{tabular}
\par\smallskip\begin{minipage}{\textwidth}\footnotesize
Fonte: elaboração própria. Valores em mil Gg de CO\textsubscript{2}, antes do filtro. A diagonal integra $c$ e $u$, mas não as forças externas. Diferenças residuais entre somas de valores exibidos decorrem do arredondamento.
\end{minipage}
\end{table}



\subsection{Seleção das relações e limites de cobertura}

O backbone de referência preserva 575 arestas, ou 13,40\% das ligações de $W$, e 70,69\% do peso intersetorial (Tabela~\ref{tab:sensibilidade-backbone}). Os painéis da Figura~\ref{fig:matriz-backbone} distinguem presença e peso: nem todas as ligações retidas correspondem a grandes volumes absolutos.

\begin{table}[htbp]
\centering
\caption{Sensibilidade da seleção e da cobertura ao parâmetro do filtro.}
\label{tab:sensibilidade-backbone}
\small
\begin{tabular}{@{}crrrr@{}}
\toprule
$\alpha$ & Arestas & \% das arestas de $W$ & Peso retido & Cobertura (\%) \\
\midrule
0,01 & 338 & 7,88 & 166,88 & 57,39 \\
0,05 & 575 & 13,40 & 205,55 & 70,69 \\
0,10 & 755 & 17,60 & 224,92 & 77,35 \\
\bottomrule
\end{tabular}
\par\smallskip\begin{minipage}{\textwidth}\footnotesize
Fonte: elaboração própria. Peso em mil Gg de CO\textsubscript{2}. Cobertura relativa ao peso de $W$, sem diagonal. Todos os casos mantêm 67 nós, um isolado e um componente fracamente conectado de 66 atividades. O filtro é aplicado à mesma rede completa, com regra OU e sem atualização dos perfis.
\end{minipage}
\end{table}

A cobertura agregada oculta perdas locais. Em $\alpha=0{,}05$, as menores coberturas são 32,96\% nas saídas e 33,79\% nas entradas. Portanto, preservar parcela expressiva do peso total não assegura representação igualmente completa de cada atividade e direção.

A mudança de $\alpha$ altera as posições individuais por grau, embora o componente de 66 atividades permaneça nos três casos. A conectividade agregada não demonstra estabilidade das posições setoriais. Os valores detalhados das mudanças e a decomposição da seleção por extremidade constam do Anexo~\ref{anexo:diagnosticos}; $0{,}05$ permanece como referência, sem ser considerado um parâmetro ótimo.



\clearpage
\newgeometry{margin=1.5cm}
\begin{landscape}
\begin{figure}[p]
\centering
\includegraphics[width=\linewidth,height=0.80\textheight,keepaspectratio]{figuras/indicadores.pdf}
\caption{Volume, concentração, conectividade selecionada e cobertura em atividades de maior volume externo.}
\label{fig:indicadores-resultados}
\begin{minipage}{\linewidth}\footnotesize
Fonte: elaboração própria. Linhas: união das seis maiores forças de saída e das seis maiores forças de entrada de $W$; não representam todas as atividades. Azul: saída; laranja: entrada. Volume e número efetivo são calculados em $W$ completa. Grau e cobertura referem-se a $\alpha=0{,}05$; os intervalos do grau abrangem $\alpha=0{,}01$--$0{,}10$ e não são intervalos de confiança. A diagonal é excluída.
\end{minipage}
\end{figure}
\end{landscape}
\restoregeometry

\clearpage
\newgeometry{margin=1.5cm}
\begin{landscape}
\begin{figure}[p]
\centering
\includegraphics[width=\linewidth,height=0.82\textheight,keepaspectratio]{figuras/matriz_presenca.pdf}
\caption{Matriz de adjacência do backbone com $\alpha=0{,}05$: (a) presença das arestas.}
\label{fig:matriz-backbone}
\begin{minipage}{\linewidth}\footnotesize
Fonte: elaboração própria. Linha: origem emissora; coluna: atividade do produto final. Ordem setorial do IBGE nos dois eixos. Azul identifica as 575 ligações retidas; cinza indica ausência no backbone, inclusive a diagonal. Ausência no backbone não implica valor nulo em $F$ ou $W$. Todas as 67 atividades permanecem nos eixos.
\end{minipage}
\end{figure}
\clearpage
\begin{figure}[p]
\ContinuedFloat
\centering
\includegraphics[width=\linewidth,height=0.82\textheight,keepaspectratio]{figuras/matriz_peso.pdf}
\caption{Matriz de adjacência do backbone com $\alpha=0{,}05$: (b) pesos das mesmas arestas.}
\begin{minipage}{\linewidth}\footnotesize
Fonte: elaboração própria. Mesma ordem e mesma seleção do painel (a). A cor representa pesos originais em Gg de CO\textsubscript{2}, em escala logarítmica, sem normalização por linha ou coluna. Cinza identifica células sem aresta retida. A cor mede volume absoluto; a concentração é descrita pelos números efetivos em $W$ completa.
\end{minipage}
\end{figure}
\end{landscape}
\restoregeometry

\subsection{Energia e construção: perspectivas de uma mesma relação}

Energia e construção ilustram uma origem emissora associada a muitos destinos finais e um destino que reúne emissões de muitas origens. A Figura~\ref{fig:rede-focos} apresenta duas vistas do mesmo backbone, com destaque para cada atividade. A interpretação das relações considera o motivo da retenção e a cobertura dos volumes, apresentados na Tabela~\ref{tab:motivos-casos} e nos painéis do Anexo~\ref{anexo:paineis}.

\begin{table}[htbp]
\centering
\caption{Motivos de retenção e cobertura nos recortes de energia e construção.}
\label{tab:motivos-casos}
\small
\begin{tabular}{@{}lrrrrr@{}}
\toprule
Recorte & Arestas & Só saída & Só entrada & Ambos & Cobertura (\%) \\
\midrule
Energia: entradas & 7 & 1 & 5 & 1 & 42,1 \\
Energia: saídas & 62 & 0 & 56 & 6 & 99,2 \\
Construção: entradas & 45 & 41 & 0 & 4 & 91,4 \\
Construção: saídas & 5 & 4 & 1 & 0 & 57,0 \\
\bottomrule
\end{tabular}
\par\smallskip\begin{minipage}{\textwidth}\footnotesize
Fonte: elaboração própria. $\alpha=0{,}05$. Saída e entrada nas colunas de motivo identificam o teste na origem e no destino da aresta, independentemente da atividade em destaque. Cobertura: peso retido dividido pela força da mesma atividade e direção em $W$. Os recortes compartilham arestas e não devem ser somados para obter totais da rede.
\end{minipage}
\end{table}

Em energia, 56 das 62 saídas retidas passam somente pelo teste de entrada no destino. Energia tem participação destacada na composição externa de muitas demandas finais, embora cada destino possa representar parcela pequena das emissões distribuídas pela origem. Entre os maiores destinos estão comércio, administração pública, construção e atividades alimentares.

Em construção, 41 das 45 entradas passam somente pelo teste de saída na origem. Construção representa parcela destacada das atribuições externas dessas origens, mesmo quando a contribuição individual é pequena no total recebido pelo destino. As entradas selecionadas pelos dois testes vêm de produtos de minerais não metálicos, energia, água e resíduos e comércio. A relação energia $\rightarrow$ construção, comum aos dois recortes, pesa aproximadamente 2\,263 Gg: 7,2\% das saídas externas de energia e 9,6\% das entradas externas de construção.

As saídas de energia e as entradas de construção têm as maiores coberturas dos respectivos recortes (Tabela~\ref{tab:motivos-casos}). Nas direções inversas, as listas curtas não autorizam desprezar as relações omitidas em conjunto. Tampouco a diferença entre 62 saídas de energia e 45 entradas de construção indica concentração distinta: na rede completa, os respectivos números efetivos são próximos, ambos de 20,8 contrapartes.



\clearpage
\begin{figure}[p]
\centering
\includegraphics[width=\textwidth,height=0.76\textheight,keepaspectratio]{figuras/rede_foco_3500.pdf}
\caption{Duas perspectivas do mesmo backbone: (a) destaque para energia (3500).}
\label{fig:rede-focos}
\begin{minipage}{\textwidth}\footnotesize
Fonte: elaboração própria. Cada vista mantém 67 nós e 575 arestas em $\alpha=0{,}05$. Azul: saídas do destaque; laranja: entradas; cinza: demais arestas. Energia tem 62 saídas e sete entradas, com cobertura de 99,2\% e 42,1\% do volume externo de cada direção. Área dos nós uniforme; espessura cresce com a raiz quadrada do peso. Posição definida pelo sinal de $u-c$ e pela ordem IBGE: 33 nós com $c>u$ à esquerda, 33 com $u>c$ à direita e o isolado abaixo. A disposição não resulta de centralidade ou comunidade.
\end{minipage}
\end{figure}
\clearpage
\begin{figure}[p]
\ContinuedFloat
\centering
\includegraphics[width=\textwidth,height=0.76\textheight,keepaspectratio]{figuras/rede_foco_4180.pdf}
\caption{Duas perspectivas do mesmo backbone: (b) destaque para construção (4180).}
\begin{minipage}{\textwidth}\footnotesize
Fonte: elaboração própria. Mesmos nós, arestas, posições e escala de espessura do painel (a). Construção tem cinco saídas e 45 entradas, com cobertura de 57,0\% e 91,4\%. Cores representam direção em relação ao destaque, não motivo de retenção. As duas vistas preservam 70,69\% do peso intersetorial de $W$. Cruzamentos e proximidade não representam trajetórias ou distâncias econômicas; relações individuais podem ser consultadas na matriz e nos painéis do anexo.
\end{minipage}
\end{figure}
\clearpage

\section{Discussão e implicações para a coordenação da mitigação}

A identificação das contrapartes esclarece o conteúdo econômico das diferenças entre produção e demanda. No cenário analisado, energia distribui emissões para muitos produtos finais, enquanto construção reúne emissões de diversas origens. O contraste mostra por que o saldo de uma atividade precisa ser interpretado junto dos volumes e das contrapartes que o compõem. Essa leitura relacional complementa as comparações de atribuição entre setores, regiões e cadeias produtivas presentes na literatura brasileira \parencite{imori2015,sanguinetAzzoni2024,montoyaEtAl2026}.

Para a coordenação da mitigação, origens compartilhadas por várias demandas e demandas associadas a várias origens indicam conjuntos de atividades a investigar conjuntamente. A geração de energia pode ser examinada junto das atividades cuja demanda final está associada às suas emissões; a construção, junto das origens que compõem sua atribuição pela demanda. Intervenções nos processos produtivos e no uso de energia e materiais ao longo das cadeias constituem frentes complementares \parencite{rissman2020}. A atribuição contábil, contudo, não determina quem deve arcar com os esforços nem identifica quais instrumentos seriam eficazes ou menos custosos.

A importância local de uma ligação também não equivale à sua contribuição para o total de emissões. O filtro permite identificar relações destacadas para uma origem ou um destino, mas as coberturas desiguais impedem tratar o backbone como uma representação igualmente completa de todas as atividades. A interpretação de oportunidades de coordenação precisa considerar os pesos e as relações omitidas, além das ligações selecionadas. A sensibilidade dos graus reforça que a presença de muitas ligações no desenho não estabelece uma hierarquia de prioridades.

As interpretações ambientais permanecem condicionadas à base utilizada. A interpolação, a cobertura não reconciliada e a diferença de preços limitam a aplicação dos contrastes setoriais à economia brasileira observada. A consistência contábil e a avaliação do filtro verificam aspectos distintos da validade das emissões físicas. Antes de sustentar conclusões ambientais definitivas, é necessário reconciliar a fronteira, a alocação setorial e os denominadores monetários dos coeficientes ou substituí-los por uma extensão compatível com 2015.

\section{Conclusões}

No cenário analisado, as diferenças entre as contas pela produção e pela demanda são compostas por relações de atribuição com volumes e concentrações distintos. Energia e construção ilustram, respectivamente, uma origem associada a muitos produtos finais e um destino que reúne emissões de diversas origens. Identificar essas contrapartes permite compreender a composição das contas além dos totais e saldos setoriais.

O backbone destaca contrapartes, mas suas perdas de cobertura e a sensibilidade dos graus exigem leitura conjunta com a rede completa. Os resultados informam questões de coordenação da mitigação no cenário, sem determinar responsabilidades ou prioridades de política. Conclusões empíricas definitivas sobre o Brasil dependem da validação ou substituição da extensão ambiental e da reaplicação da mesma especificação.

{\renewcommand*{\bibfont}{\small}
\printbibliography[title={Referências}]
}
\clearpage
\appendix
\section{Painéis completos de energia e construção}
\label{anexo:paineis}

\begin{figure}[!ht]
\centering
\includegraphics[width=\textwidth,height=0.78\textheight,keepaspectratio]{figuras/anexo_energia_1.pdf}
\caption{Energia (3500): todas as sete entradas e as saídas 1--31 de 62.}
\label{fig:anexo-energia}
\begin{minipage}{\textwidth}\footnotesize
Fonte: elaboração própria. Recortes de $\alpha=0{,}05$, com pesos originais em Gg de CO\textsubscript{2}, uma casa decimal e escala linear comum. A cobertura é o peso retido sobre a força da atividade e direção em $W$, sem diagonal: 42,1\% nas entradas e 99,2\% nas 62 saídas. As demais saídas aparecem na continuação, sem novo filtro. Azul: só saída na origem; laranja: só entrada no destino; roxo: ambos os testes.
\end{minipage}
\end{figure}
\clearpage
\begin{figure}[p]
\ContinuedFloat
\centering
\includegraphics[width=\textwidth,height=0.85\textheight,keepaspectratio]{figuras/anexo_energia_2.pdf}
\caption{Energia (3500), continuação: saídas 32--62 de 62.}
\begin{minipage}{\textwidth}\footnotesize
Fonte: elaboração própria. Mesma seleção, ordem, escala e cores do painel anterior. A cobertura de 99,2\% se refere às 62 saídas, não apenas a esta página. Azul: só saída na origem; laranja: só entrada no destino; roxo: ambos os testes.
\end{minipage}
\end{figure}
\clearpage
\begin{figure}[p]
\centering
\includegraphics[width=\textwidth,height=0.85\textheight,keepaspectratio]{figuras/anexo_construcao.pdf}
\caption{Construção (4180): todas as 45 entradas e cinco saídas.}
\label{fig:anexo-construcao}
\begin{minipage}{\textwidth}\footnotesize
Fonte: elaboração própria. Coberturas de 91,4\% nas entradas e 57,0\% nas saídas. Mesma escala de pesos dos painéis de energia. Azul: só saída na origem; laranja: só entrada no destino; roxo: ambos os testes. Arestas compartilhadas não são duplicadas na rede.
\end{minipage}
\end{figure}
\clearpage

\section{Diagnósticos complementares da rede e da seleção}
\label{anexo:diagnosticos}

Na rede completa, os 66 perfis definidos têm medianas de 11,50 destinos efetivos e 13,93 origens efetivas, abaixo das 65 contrapartes existentes em cada direção. No backbone com $\alpha=0{,}05$, as medianas de cobertura são 59,32\% nas saídas e 55,69\% nas entradas.

A Tabela~\ref{tab:postos-complementares} detalha as mudanças de posição por grau em relação a $\alpha=0{,}05$. Os postos são calculados em ordem decrescente entre os 67 nós, com posto médio nos empates; reporta-se a mediana da mudança absoluta.

\begin{table}[htbp]
\centering
\caption{Mudança mediana de posição por grau em relação a $\alpha=0{,}05$.}
\label{tab:postos-complementares}
\begin{tabular}{@{}crr@{}}
\toprule
$\alpha$ & Grau de saída & Grau de entrada \\
\midrule
0,01 & 11 & 7 \\
0,10 & 4,5 & 4 \\
\bottomrule
\end{tabular}
\par\smallskip\begin{minipage}{\textwidth}\footnotesize
Fonte: elaboração própria. Valores em posições; as diferenças são calculadas a partir dos postos não arredondados, incluindo empates.
\end{minipage}
\end{table}

Dos 575 vínculos retidos em $\alpha=0{,}05$, 232 passam somente pelo teste de saída na origem, 265 somente pelo de entrada no destino e 78 pelos dois. Nenhum depende da convenção de grau um. A seleção por apenas uma extremidade mostra que uma relação pode destacar-se em um perfil sem ter a mesma importância relativa no outro.

Nos recortes da Tabela~\ref{tab:motivos-casos}, o backbone omite 57,9\% das entradas externas de energia e 43,0\% das saídas externas de construção. Esses complementos das coberturas evidenciam o volume omitido nas direções com menos ligações selecionadas.

\end{document}
